\documentclass[11pt,a4paper]{article}
\usepackage[utf8]{inputenc}
\usepackage[T1]{fontenc}
\usepackage[french]{babel}
\usepackage{amsmath,amssymb,amsthm,geometry,hyperref}
\geometry{margin=2.4cm}
\newtheorem{proposition}{Proposition}
\title{Cubature caractéristique à moments compatibles\\Note théorique du prototype SLFE GPU}
\author{Note de recherche expérimentale}
\date{15 novembre 2026}
\begin{document}
\maketitle

\section{Problème et hypothèses}
Soit $F:\Omega_s\to\Omega_t$ un difféomorphisme préservant le volume,
sans échange aux frontières, éventuellement périodiques. Les espaces éléments
finis $V_s,V_t$ contiennent la constante. Leurs matrices de masse exactes,
symétriques définies positives, sont notées $M_s,M_t$. Les matrices
d'évaluation $E_s,E_t$ sont définies sur des points $a_\ell$ et leurs images
$F(a_\ell)$. Une même matrice diagonale $W=\operatorname{diag}(w_\ell)$
sert à construire le transfert
\[
 T=M_t^{-1}E_t^TWE_s.
\]
Il s'agit de la réalisation par cubature du transfert faible conservatif.
Les poids concernent la cubature, et ne sont pas les coefficients du transfert
dans une base modale ou nodale.

\begin{proposition}
Si $w_\ell\geq0$ et
\[
 E_s^TWE_s=M_s,\qquad E_t^TWE_t=M_t,
\]
alors $T$ conserve la masse, reproduit la constante et est contractant pour
les normes discrètes $L^2$ définies par les matrices de masse.
\end{proposition}
\begin{proof}
Soient $e_s,e_t$ les coefficients représentant la constante, de sorte que
$E_se_s=E_te_t=\mathbf1$. Alors
\[
 M_tTe_s=E_t^TW\mathbf1=E_t^TWE_te_t=M_te_t,
\]
et
\[
 e_t^TM_tTc=e_t^TE_t^TWE_sc
 =e_s^TE_s^TWE_sc=e_s^TM_sc.
\]
Enfin, les matrices $U_s=W^{1/2}E_sM_s^{-1/2}$ et
$U_t=W^{1/2}E_tM_t^{-1/2}$ ont des colonnes orthonormales. Par conséquent,
\[
 M_t^{1/2}TM_s^{-1/2}=U_t^TU_s,
 \qquad \|Tc\|_{M_t}\leq\|c\|_{M_s}.
\]
\end{proof}

Cette proposition ne garantit pas la positivité du champ reconstruit.
L'orthogonalité, la conservation, la monotonie et la positivité sont des
propriétés distinctes. Elle ne démontre pas non plus l'existence de poids
positifs sur un ensemble de particules fixé.

Pour des éléments affines DG $P_p$, chaque produit de fonctions de base
appartient à $P_{2p}$. Les égalités de Gram s'obtiennent donc en imposant
les moments jusqu'au degré $2p$ séparément dans chaque élément de départ
et d'arrivée. La représentation modale orthonormale donne une matrice de
masse locale égale à l'aire multipliée par l'identité.

\section{Précision et défauts de cubature}
La contraction ne suffit pas à établir un ordre de convergence. Il faut
contrôler séparément l'erreur de consistance du couplage croisé
$E_t^TWE_s$, celle des trajectoires et leur accumulation temporelle.
Si $F$ est affine et si un polynôme transporté est représentable dans les
deux espaces, les moments d'arrivée jusqu'au degré $2p$ suffisent à sa
reproduction, sous réserve de compatibilité avec les frontières.
Une translation périodique d'un champ lisse demande une véritable étude
de raffinement, et ne se réduit pas à ce test polynomial.

Si les défauts de Gram normalisés vérifient
\[
 \|M_i^{-1/2}(E_i^TWE_i-M_i)M_i^{-1/2}\|_2\leq\varepsilon_i,
\]
alors, pour des poids non négatifs,
\[
 \|M_t^{1/2}TM_s^{-1/2}\|_2
 \leq\sqrt{(1+\varepsilon_s)(1+\varepsilon_t)}.
\]
Un petit résidu du solveur de moments ne doit donc pas être présenté comme
une preuve de contraction exacte.

\section{Extension compressible}
Pour $J_\ell=\det DF(a_\ell)>0$, une construction compatible utilise
\[
 E_s^TWE_s=M_s,\qquad E_t^T\operatorname{diag}(J_\ell)WE_t=M_t,
 \qquad T=M_t^{-1}E_t^TWE_s.
\]
La masse reste conservée. Le champ constant n'est généralement pas invariant
pour une densité conservative. Le même argument donne
\[
 \|Tc\|_{M_t}\leq (\min_\ell J_\ell)^{-1/2}\|c\|_{M_s}.
\]
La précision du Jacobien et la fidélité des points transportés sont ici des
conditions supplémentaires. Une masse numérique conservée ne suffit pas à
valider l'équation compressible.

\section{Construction expérimentale et limites}
Le prototype ajuste les poids d'une quadrature de départ en minimisant
\[
 \frac12\sum_\ell\frac{(w_\ell-w_\ell^0)^2}{w_\ell^0}
 \quad\text{sous les contraintes de moments des deux maillages.}
\]
La variante signée est un problème de moindres carrés. La variante positive
expérimentale supprime les sites obtenant des poids négatifs puis recalcule
les poids. Cette stratégie restreinte peut échouer même si une autre
configuration active est faisable. Elle ne constitue pas un solveur complet
de programmation quadratique ni un certificat d'infaisabilité.

Les systèmes peuvent être très mal conditionnés : les contraintes de départ
et d'arrivée deviennent presque dépendantes pour certaines combinaisons
polynomiales. Une SVD GPU dense sert uniquement de diagnostic sur de petits
systèmes. Une réalisation utile sur de grands maillages demande un traitement
de cette structure, un préconditionnement efficace ou une construction locale.

Une variante par projections alternées de Dykstra est également disponible
pour le diagnostic du problème quadratique positif. Sa limite d'itérations
doit être distinguée de la faisabilité mathématique. Les expériences de
compression montrent que l'enrichissement du support peut rétablir les
contraintes alors que la recherche sur un support plus pauvre échoue.

\section{Obstruction du support à petit déplacement}
\begin{proposition}
Considérons une translation $F(a)=a+\delta$, $\delta\ne0$, et un même
maillage de départ et d'arrivée. Supposons qu'une cellule $K$ n'échange
aucune particule avec ses voisines : tous ses points restent dans $K$ et
aucun point extérieur n'y arrive. Il est impossible d'imposer simultanément
les moments de degré zéro et un exacts au départ et à l'arrivée dans $K$,
même avec des poids signés.
\end{proposition}
\begin{proof}
Dans une carte locale, les mêmes indices $I_K$ interviennent des deux côtés.
Les contraintes imposent
\[
 \sum_{\ell\in I_K}w_\ell=|K|,\qquad
 \sum_{\ell\in I_K}w_\ell a_\ell=\int_K x\,dx,
 \qquad
 \sum_{\ell\in I_K}w_\ell(a_\ell+\delta)=\int_K x\,dx.
\]
La soustraction des deux dernières égalités donne $|K|\delta=0$, une
contradiction. L'argument ne fait aucune hypothèse de signe sur les poids.
\end{proof}
Pour un support fini strictement intérieur aux cellules, un déplacement
non nul suffisamment petit ne fait traverser aucune interface. Il existe
donc un régime où la compatibilité est impossible pour tout $p\geq1$.
Ce résultat concerne un support fixe. Il ne condamne pas une construction
qui adapte les points aux interfaces et au déplacement, mais démontre
qu'augmenter une fois pour toutes le degré de quadrature ne résout pas le
problème lorsque le pas tend vers zéro. Les compteurs de cellules isolées
de la campagne permettent de détecter cette obstruction géométrique,
distincte de l'échec d'un solveur sur un support potentiellement faisable.

L'intérêt potentiel est de conserver des trajectoires non linéaires exactes
sans construire leurs cellules images courbes. Les gains GPU, la convergence,
la robustesse à la déformation et l'originalité par rapport aux travaux de
remaillage conservatif par optimisation restent à établir expérimentalement
et bibliographiquement.
\end{document}
